\section{Introduction}\label{sec:intro}
Isoparametric elements lose accuracy on distorted meshes because the parametric shape
functions, composed with a non-affine element map, no longer span the complete polynomials
in Cartesian coordinates~\cite{Lee1993effects,Arnold2002Approximation}. The unsymmetric
finite element method of \citet{Rajendran2003novel} restores completeness by a
Petrov--Galerkin formulation: the test functions remain parametric, the trial functions are
\emph{metric}, i.e.\ polynomials in physical coordinates that interpolate the nodal values.
With the Jacobian determinant cancelled in the stiffness integral, the resulting elements
reproduce every displacement field of total degree $\le p$ exactly on arbitrarily distorted
meshes, also when $\det\bm J$ changes sign (Jacobian-unlimited elements, conditions C1--C3 of
\citet{Rajendran2010technique}). The approach has been developed for a wide range of element
types and applications, e.g.\
\cite{Ooi2004hexahedron,Ooi2007extension,Cen2012shape,Cen2015unsymmetric,Xie2016Modified,Shang2018airy,Huang2022Unsymmetric,Pfefferkorn2021mesh,Pfefferkorn2023hourglassing,Zaehringer2026On};
two known difficulties are the frame dependence of Cartesian metric functions, removed by
skew (corotational) coordinates~\cite{Ooi2008Remedies,Xie2016Modified}, and the
``interpolation failure'' of the metric functions on strongly distorted
elements~\cite{Ooi2008Remedies,Rajendran2010technique}.

For high-order (spectral) elements, the metric trial functions of degree $p$ in each
coordinate are obtained element by element from a generalised Vandermonde system with
$(p+1)^2$ unknowns. Two different properties of this construction degrade with $p$ and with
the element distortion, and they are often not distinguished:
\begin{enumerate}
\item the condition number $\kappa(\bm V)$ of the interpolation matrix, which controls the
  round-off in the computed shape functions and depends on the polynomial basis used to set up
  the system;
\item the Lebesgue constant $\Lambda$ of the interpolation in the trial space, which controls the
  approximation and the amplification of nodal perturbations; it is a property of the space and
  the node set and does not depend on the basis.
\end{enumerate}
A companion study of the thick cylinder of \citet[Sec.~5.9]{Rajendran2010technique} with
$p=2,\dots,8$ (\cref{sec:cylinder}) shows the practical consequence: on meshes whose relative
distortion does not decrease under refinement, the unsymmetric spectral element with metric
functions $Q_p$ in skew coordinates converges optimally only for moderate distortion; beyond, its
convergence is erratic, and the loss correlates with Lebesgue constants of $10^{3}$--$10^{8}$
rather than with the sign of $\det\bm J$.

This paper proposes one remedy for each property.
\begin{enumerate}
\item \emph{A well-conditioned basis} for the metric interpolation problem
  (\cref{sec:basis}): a tensor-product Lagrange basis on the Gauss--Lobatto--Legendre (GLL)
  points of the bounding box of the skew node coordinates. We prove that the trial space and its
  Lebesgue constant are basis invariant, that the new basis makes $\bm V$ a permutation matrix on
  affine elements, and that on distorted elements $\kappa_1(\bm V)$ equals the product of the
  reference Lebesgue function at the scaled nodes and the Lebesgue function of the metric
  interpolation at the reference points, i.e.\ the basis removes all ill-conditioning that is not
  inherent in the interpolation problem. It does \emph{not} reduce the Lebesgue constant.
\item \emph{A hybrid trial space} (\cref{sec:hybrid}) that uses metric polynomials only where
  completeness requires them, for $P_p(\bx)$, and parametric functions for the remaining
  $\tfrac12p(p+1)$ functions. We establish completeness, frame invariance and the reduction to
  the conventional spectral element on bilinear elements, derive a factorisation of the
  interpolation matrix that characterises unisolvence, and show that the hybrid interpolant is
  the GLL interpolant plus the interpolation error of a polynomial, which bounds its Lebesgue
  constant by the GLL value plus a term that vanishes on bilinear elements.
\end{enumerate}
\Cref{sec:numerics} supports the claims by element-level studies up to $p=16$, mesh-level
studies (smooth fields, eigenvalues), Jacobian-unlimited patch tests on the meshes of
\citet{Rajendran2010technique}, and the thick-cylinder study with $1512$ solutions.
\Cref{sec:conclusions} summarises the advantages and the limitations.
